\documentclass{beamer}
\usepackage[utf8]{inputenc}
\usepackage[T1]{fontenc}
\usepackage{amsmath}

\usetheme{Madrid}
\usecolortheme{whale}

\title{Aprendizado de Máquina: Uma Introdução}
\subtitle{Conceitos Fundamentais}
\author{Seu Nome}
\date{\today}
\institute{Sua Instituição}

\begin{document}

\begin{frame}
\titlepage
\end{frame}

\begin{frame}{Agenda}
\tableofcontents
\end{frame}

\section{O que é Aprendizado de Máquina?}

\begin{frame}{Definição}
  Aprendizado de Máquina é um subcampo da IA que permite que sistemas
  \textbf{aprendam a partir de dados} sem serem explicitamente programados.

  \bigskip
  Três paradigmas principais:
  \begin{itemize}
    \item \textbf{Supervisionado} — aprende com pares (entrada, saída)
    \item \textbf{Não supervisionado} — encontra padrões sem rótulos
    \item \textbf{Por reforço} — aprende por tentativa e erro
  \end{itemize}
\end{frame}

\section{Regressão Linear}

\begin{frame}{Regressão Linear}
  Dado um conjunto $\{(x_i, y_i)\}_{i=1}^n$, queremos encontrar:

  \[
  \hat{y} = \theta_0 + \theta_1 x
  \]

  Minimizando o erro quadrático médio:

  \[
  J(\theta) = \frac{1}{2n} \sum_{i=1}^{n} \left(\hat{y}_i - y_i\right)^2
  \]
\end{frame}

\begin{frame}{Gradiente Descendente}
  Atualização dos parâmetros:

  \[
  \theta_j \leftarrow \theta_j - \alpha \frac{\partial J}{\partial \theta_j}
  \]

  onde $\alpha$ é a taxa de aprendizado.

  \begin{block}{Intuição}
    Seguimos o gradiente negativo — a direção de maior descida.
  \end{block}
\end{frame}

\section{Conclusão}

\begin{frame}{Próximos Passos}
  \begin{enumerate}
    \item Estudar regularização (L1, L2)
    \item Explorar redes neurais
    \item Praticar com datasets reais (Kaggle, UCI)
  \end{enumerate}

  \bigskip
  \centering
  \textit{``A melhor forma de aprender é fazendo.''}
\end{frame}

\end{document}
